\noindent\textbf{ON THE HANOVER TRIANGULATION.}\hfill347
\begin{center}\textsc{Gauss} to \textsc{Schumacher}.\quad Göttingen, 10 September 1818.\end{center}

I hasten to tell you that our minister, \textsc{Arnswaldt}, has charged me
with carrying out in Lüneburg the measurements needed to link a Hanover
triangulation with yours and making the necessary
arrangements with you. He also gives me reason to hope that
the extension itself will probably come to pass before long, and
I am glad that it can now be secured by the operations to be carried out
in Lüneburg.\ .\ .\ .\ .\ .\ .

\begin{center}\textsc{Gauss} to \textsc{Schumacher}.\quad Göttingen, 20 May 1820.\end{center}
\begingroup
\makeatletter
\renewcommand{\@makefnmark}{\hbox{\normalfont[\textsuperscript{*})]}}
\long\def\@makefntext#1{\parindent1.5em\noindent[\textsuperscript{*})\ #1]}
\makeatother
\par\noindent\hspace*{\parindent}.\ .\ .\ .\ .\ .\ In this uncertainty\footnote{About \textsc{Schumacher}'s whereabouts.} I am addressing this letter to Copenhagen,
and I earnestly hope that it reaches you, and reaches you \textls[120]{soon}. Its
purpose is to bring you the news that, following a letter
from Count \textsc{von Münster} in London in reply to my letter of a year ago,
\makeatletter
\renewcommand{\@makefnmark}{\hbox{\normalfont[\textsuperscript{**})]}}
\long\def\@makefntext#1{\parindent1.5em\noindent[\textsuperscript{**})\ #1]}
\makeatother
which you kindly saw to for me\footnote{Printed in Vol. IV, pp.~482/483.},
\begin{quote}
»the King has approved the extension of the arc measurement through the Kingdom
of Hanover«.
\end{quote}
\par\medskip
.\ .\ .\ .\ .\ .
\endgroup

\begin{center}\textsc{Gauss} to \textsc{Schumacher}.\quad Göttingen, 4 March 1821.\end{center}

.\ .\ .\ .\ .\ .\ I had originally intended to start from Hamburg and thus
to measure from north to south. Unfortunately, since I have been held up so badly by the
reluctance to write and the unreliability of all the instrument makers I have dealings
with, and since I still have none of the necessary
equipment in hand, the difficulties arising from this
would be much greater if I followed that plan. This and
\textls[120]{several other important} reasons compel and induce me
to adopt the opposite plan, measuring from south to north.\ .\ .\ .\ .\ .\ .
\par
\makeatletter
\insert\footins{\reset@font\footnotesize\hsize\columnwidth\parindent0pt
\smallskip\hbox to\hsize{\hfil44\textsuperscript{*}}\par}
\makeatother

\clearpage
\noindent348\hfill\textbf{CORRESPONDENCE WITH SCHUMACHER.}
\begin{center}\textsc{Gauss} to \textsc{Schumacher}.\quad Wulfsode, 18 September 1822.\end{center}

.\ .\ .\ .\ .\ .\ As far as one can judge from a preliminary inspection of the terrain,
it would not be impossible to measure \textls[120]{the entire line} from Breithorn to Eschede
[Scharnhorst] (11220 metres long) directly. What a splendid
baseline this would be!\ .\ .\ .\ .\ .\ .

\begin{center}\textsc{Gauss} to \textsc{Schumacher}.\quad [Göttingen, January 1825.]\end{center}

.\ .\ .\ .\ .\ .\ If I count all the larger and smaller cuttings through the woods from the years
1821---1824, from those where perhaps a dozen trees
were felled to the largest, there may have been about 16 or 17 cuttings.
The very largest in extent was in the Becklinger
Holz, near the road from Bergen to Soltau.\ .\ .\ .\ .\ .\ .

\begin{center}\textsc{Gauss} to \textsc{Schumacher}.\quad Dangast, 1 hour from Varel, 20 June 1825.\end{center}

.\ .\ .\ .\ .\ .\ My own practice [with the heliotrope] is to draw a circle through the tips
of the three feet, taking its centre as the target point.
On my two newest heliotropes the centre itself is also marked by a
spike, which is very convenient.\ .\ .\ .\ .\ .\ .

\begin{center}\textsc{Gauss} to \textsc{Schumacher}.\quad Göttingen, 14 January 1829.\end{center}

Since we still know little about whether the figure of the Earth departs from the figure
of the mean ellipsoid in longer or shorter undulations,
I am in any case of the opinion that it is best always to take the mean
ellipsoid as the basis.

If, however, I am unwilling to reject outright the occasional use of an
osculating ellipsoid as the basis, I can regard it as admissible only
where there are means of determining an osculating ellipsoid,
i.e. one whose curvature, \textls[120]{both} in the direction of the meridian
and in the direction perpendicular to it, comes as close as possible to
that of the actual figure. But this can be achieved only by combining \textls[120]{measurements}
in both directions (arc measurements in latitude and longitude).\ .\ .\ .\ .\ .\ .

\clearpage
\noindent\textbf{ON THE HANOVER TRIANGULATION.}\hfill349
\par\medskip
\begingroup
\makeatletter
\renewcommand{\@makefnmark}{\hbox{\normalfont\textsuperscript{*})}}
\makeatother
If a meridian is not actually elliptical, one can indeed
calculate an ellipse that fits two portions of an arc, and
for all I care one may call this an osculating ellipse. But
rotating this ellipse about \textls[120]{its} minor axis\footnote{\hspace{.25em}All that can be said of it is that it is parallel to the Earth's axis without coinciding with it. Indeed, it will generally be far from it.} produces no surface that
can properly be called an osculating ellipsoid. In other words, there is
no relation at all between the actual value of a degree of longitude and the degree of longitude
calculated on the [ellipsoid obtained] by rotating that osculating ellipse.
Anyone who confuses the two is liable to make
the greatest errors.\ .\ .\ .\ .\ .\ .
\endgroup
\begin{center}\rule{.26\linewidth}{.4pt}\end{center}
